% s03-2-1-fusability.tex —— 3.2.1 小节：可熔性（印刷页 82–86）

\section{Arkhipov–Bezukavnikov 等价}
\label{sec:ab}

\subsection{可熔性}
\label{sec:fusability}

本小节讨论\emph{可熔的}（fusable）因子化 $\Rep(\check G)$-模范畴这一概念。该概念将在
Bogdanova、Chen 与 Gaitsgory 即将发表的文章 \cite{BCG26} 中仔细定义。

考虑
\[
\LS^{\mer}_{\check G,x_0}=\LS_{\check G}(D^\times_{x_0}).
\]
于是可作对称幺半范畴 $\QCoh(\LS_{\check G}(D^\times_{x_0}))$ 上的模范畴构成的（$2$-，
即 $(\infty,2)$-）范畴，记作
\[
\QCoh(\LS_{\check G}(D^\times_{x_0}))\text{-}\ModCat.
\]
在 \cite[3.1 节]{Bog24} 中，Bogdanova 构造了（$2$-）函子
\[
\mathrm{Fact}\colon
\QCoh(\LS_{\check G}(D^\times_{x_0}))\text{-}\ModCat
\longrightarrow
\Rep(\check G)\text{-}\FactModCat_{x_0},
\]
它由
\[
\mathcal{M}\mapsto
\QCoh(\LS_{\check G}(D^\times_{\Ran_{x_0}}))
\otimes_{\QCoh(\LS_{\check G}(D^\times_{x_0}))}\mathcal{M}
\]
给出。该函子被猜测为全忠实；该文朝此猜测取得了如下进展。
\srcpage{82}

在 \cite[3.3 节]{Bog24} 中给出了限制局部系统的叠
$\LS^{\mathrm{restr}}_{\check G}(D^\times_{x_0})$ 的定义，并有映射
\[
\LS^{\mathrm{restr}}_{\check G}(D^\times_{x_0})
\longrightarrow \LS_{\check G}(D^\times_{x_0}),
\]
它经拉回定义幺半函子
\[
\QCoh(\LS_{\check G}(D^\times_{x_0}))\longrightarrow
\QCoh(\LS^{\mathrm{restr}}_{\check G}(D^\times_{x_0})),
\]
后者又经限制给出
\[
\QCoh(\LS^{\mathrm{restr}}_{\check G}(D^\times_{x_0}))\text{-}\ModCat
\longrightarrow
\QCoh(\LS_{\check G}(D^\times_{x_0}))\text{-}\ModCat.
\]
把复合
\begin{align*}
&\QCoh(\LS^{\mathrm{restr}}_{\check G}(D^\times_{x_0}))\text{-}\ModCat
\longrightarrow \QCoh(\LS_{\check G}(D^\times_{x_0}))\text{-}\ModCat\\
&\qquad\quad
\xrightarrow{\ \mathrm{Fact}\ }\Rep(\check G)\text{-}\FactModCat_{x_0}
\end{align*}
记作 $\mathrm{Fact}^{\mathrm{restr}}$。Bogdanova 在该文定理 3.8 中证明了

\begin{proposition}
\label{prop:fact-restr-ff}
$\mathrm{Fact}^{\mathrm{restr}}$ 全忠实。此外，图
\[
\begin{tikzcd}[column sep=large,row sep=large]
\QCoh(\LS^{\mathrm{restr}}_{\check G}(D^\times_{x_0}))\text{-}\ModCat
\arrow[rr,"\mathrm{Fact}^{\mathrm{restr}}"]
\arrow[dr,"\oblv"'] & &
\Rep(\check G)\text{-}\FactModCat_{x_0}
\arrow[dl,"\mathrm{cores}_{x_0}"]\\
& \DGCat
\end{tikzcd}
\]
交换。\yizhu{原文图中出现记号 $\mathrm{cores}_{x_0}$，但全文未给出它的定义；
按上下文推断，它应指 $x_0$ 处的余限制函子。}
{译者按（上下文推断，无外部依据）}
\end{proposition}

下面给出可熔性的定义。

\begin{definition}
\label{def:fusable}
因子化 $\Rep(\check G)$-模范畴 $\mathcal{M}$ 称为\emph{可熔的}，若它落在函子
\[
\mathrm{Fact}^{\mathrm{restr}}\colon
\QCoh(\LS^{\mathrm{restr}}_{\check G}(D^\times_{x_0}))\text{-}\ModCat
\longrightarrow \Rep(\check G)\text{-}\FactModCat_{x_0}
\]
的本质像中。
\end{definition}
\srcpage{83}

把 $\Rep(\check G)\text{-}\FactModCat_{x_0}$ 中可熔对象构成的全子范畴记作
\[
\Rep(\check G)\text{-}\FactModCat^{\mathrm{fus}}_{x_0}
\subset \Rep(\check G)\text{-}\FactModCat_{x_0}.
\]

\begin{remark}
\label{rem:fusable-combinatorial}
\cite{BCG26} 将给出可熔性的组合刻画，使得在实践中可以判定特定的因子化
$\Rep(\check G)$-模范畴是否可熔。现简述这一刻画：

给定因子化 $\Rep(\check G)$-模范畴 $\mathcal{M}$，可考虑它到曲线的单独一份上的限制
$\mathcal{M}_X$。记
\[
U\xrightarrow{\ j\ }X\xleftarrow{\ i\ }x_0,
\]
注意进一步限制到 $U$ 上是
\[
\mathcal{M}_U\simeq \Rep(\check G)\otimes D\text{-}\mathrm{mod}(U)\otimes \mathcal{M}_{x_0}.
\]
有函子
\[
j_{\bullet}\colon \Rep(\check G)\otimes D\text{-}\mathrm{mod}(U)\otimes \mathcal{M}_{x_0}
\simeq \Gamma(U,\mathcal{M}_U)\longrightarrow \Gamma(X,\mathcal{M}_X),
\]
以及函子
\[
i^{\bullet}\colon \Gamma(X,\mathcal{M}_X)\longrightarrow \mathrm{Pro}(\mathcal{M}_{x_0}).
\]
可把复合
\[
i^{\bullet}j_{\bullet}\colon
\Rep(\check G)\otimes D\text{-}\mathrm{mod}(U)\otimes \mathcal{M}_{x_0}
\to \mathrm{Pro}(\mathcal{M}_{x_0})
\]
限制到
\[
\Rep(\check G)\otimes \QLisse(U)\otimes \mathcal{M}_{x_0}
\subset \Rep(\check G)\otimes D\text{-}\mathrm{mod}(U)\otimes \mathcal{M}_{x_0}.
\]
\srcpage{84}
于是 $\mathcal{M}$ 可熔等价于限制
\[
\Rep(\check G)\otimes \QLisse(U)\otimes \mathcal{M}_{x_0}
\longrightarrow \mathrm{Pro}(\mathcal{M}_{x_0})
\]
的像落在 $\mathcal{M}_{x_0}$ 中而非其 pro-范畴中（再加上一些涉及曲线更高次幂的条件，
它们将出现在 \cite{BCG26} 中）。

这一概念的有用之处有两方面。首先，由于
\[
\oblv\colon \QCoh(\LS^{\mathrm{restr}}_{\check G}(D^\times_{x_0}))\text{-}\ModCat
\longrightarrow \DGCat
\]
保守，若 $F\colon \mathcal{M}_1\to \mathcal{M}_2$ 是可熔因子化
$\Rep(\check G)$-模范畴之间的因子化函子，则 $F_{x_0}$ 是等价时 $F$ 就是等价。

其次，若 $\mathcal{M}$ 是因子化 $\Rep(\check G)$-模范畴，则有因子化
$\Rep(\check G)$-子模范畴
\[
\mathcal{M}_{\QLisse}\subset \mathcal{M},
\]
其刻画如下。令 $U=X\setminus\{x_0\}$，并对每个满射 $n\to m$，记对角线为
$\Delta_{n\to m}\colon X^m\to X^n$。$\Delta_n$ 记 $X^n$ 中一切（非恒等的）对角线之并。
把 $X^n$（及其子集）经
\[
(x_1,\dots,x_n)\mapsto \{x_0,x_1,\dots,x_n\}
\]
看作映到 $\Ran_{x_0}$。由因子化有
\[
\mathcal{M}|_{U^n\setminus \Delta_n}
\simeq \bigl(\Rep(\check G)^{\otimes n}\otimes D\text{-}\mathrm{mod}(U)^{\otimes n}
\otimes \mathcal{M}_{x_0}\bigr)|_{U^n\setminus \Delta_n},
\]
\srcpage{85}
并定义 $\mathcal{M}_{\QLisse}$ 为满足
\[
\mathcal{M}_{\QLisse}|_{U^n\setminus \Delta_n}
\simeq \bigl(\Rep(\check G)^{\otimes n}\otimes \QLisse(U)^{\otimes n}
\otimes \mathcal{M}_{x_0}\bigr)|_{U^n\setminus \Delta_n}
\]
的对象构成的子范畴。

当 $\mathcal{M}$ 可熔时，要构造源为 $\mathcal{M}$ 的因子化函子，只需构造源为
$\mathcal{M}_{\QLisse}$ 的因子化函子。

下一小节将用后一事实构造对之间的态射 AB；再下一小节将用前一观察论证它是对之间的
等价。为此需要
\end{remark}

\begin{proposition}
\label{prop:whitFl-fusable}
经因子化范畴的等价 $\mathrm{CS}_G$ 看作因子化 $\Rep(\check G)$-模范畴时，
$\Whit^{!}(\Fl_G)$ 可熔。
\end{proposition}

\begin{proof}
\cite{Gai04} 研究了族 $\Fl'\to X$，即因子化模空间 $\Fl_{G,\Ran_{x_0}}$ 到曲线单独
一份上的限制。利用卷积积
\[
D\text{-}\mathrm{mod}(\Gr_G)^{L^+G}\otimes D\text{-}\mathrm{mod}(U)
\otimes D\text{-}\mathrm{mod}(\Fl_G)\longrightarrow D\text{-}\mathrm{mod}(\Fl_G),
\]
并经（逐点的）朴素 Satake 函子
$\Rep(\check G)\to D\text{-}\mathrm{mod}(\Gr_G)^{L^+G}$，得到所要的函子
$i^{\bullet}j_{\bullet}$。因此 $D\text{-}\mathrm{mod}(\Fl_G)$ 可熔；对两侧取
$\Whit^{!}$ 后，$\Whit^{!}(\Fl_G)$ 也可熔。
\end{proof}

此外，由构造有

\begin{proposition}
\label{prop:qcoh-nB-fusable}
$\QCoh(\check{\mathfrak{n}}/\check B)$ 可熔。
\end{proposition}

\begin{proof}
把 $\QCoh(\check{\mathfrak{n}}/\check B)$ 的构造展开，并与 \cite{Bog24} 中
$\mathrm{Fact}^{\mathrm{restr}}$ 的描述比较，即知
\[
\QCoh(\check{\mathfrak{n}}/\check B)
\simeq \mathrm{Fact}^{\mathrm{restr}}\bigl(\QCoh(\check{\mathfrak{n}}/\check B)_{x_0}\bigr).
\]
\end{proof}
\srcpage{86}
